\section{Transferring Files}

Move files between your laptop and the server using SCP or rsync.

\subsection{SCP}

\begin{lstlisting}
scp file.txt username@pelee.math.private.uwaterloo.ca:~
\end{lstlisting}

\subsection{RSYNC}

\begin{lstlisting}
rsync -av data/ username@pelee.math.private.uwaterloo.ca:data/
\end{lstlisting}

\subsection{Mounting a Remote Folder with SSHFS}

SSHFS makes a folder on a server appear as an ordinary folder on your Linux computer, so you can
open and edit remote files with local programs without copying them back and forth. These
instructions are for Linux (on Ubuntu or Debian: \texttt{sudo apt install sshfs}).

\begin{lstlisting}
mkdir -p ~/nibi
sshfs username@nibi.alliancecan.ca:/scratch/username ~/nibi
ls ~/nibi
fusermount -u ~/nibi        # unmount when finished
\end{lstlisting}

If you use the shortcut file of Appendix~\ref{app:ssh} (with \texttt{ControlMaster}), the
multi-factor prompt appears only for the first connection.

\begin{itemize}
\item Use a mount for browsing, editing and small files. For large transfers use
\texttt{rsync} or Globus, and do not run heavy programs on files in a mounted folder.
\item Unmount when you are done. If a mount hangs after your network changes, use
\texttt{fusermount -uz \textasciitilde/nibi}.
\end{itemize}

\subsection{Exercise}

Transfer a small text file and verify the transfer using \texttt{ls}.
